\section{Limits of the Formal Approach}
\label{sec:limits}

The Article's structural claim is bounded.
This Part identifies the limits of the claim, the objections to which the Article's argument is most vulnerable,
and the responses that are available against those objections.
The Part is offered in the spirit of candor about cross-disciplinary work:
a CS-flavored argument applied to legal-doctrinal domains should make its limits explicit,
because the asymmetry of expertise between the author's home discipline and the target audience's home discipline
is large and the failure modes of cross-disciplinary work are well-documented.

\subsection{The Schmittian objection}

Part~\ref{sec:pattern} acknowledged that for Schmitt the constitutive character of the sovereign decision
is normatively defensible under conditions of genuine emergency, not merely descriptive of how legal orders
are seen to fail, and that the structural argument this Article advances is incompatible with that
normative view rather than orthogonal to it.
This Part considers the implication for the Article's scope.

The implication is that the Article's structural claim does not bind a reader who holds the strong Schmittian
view.
A reader for whom the constitutive sovereign decision is the normatively prior fact of political order
will read the typed-rollback discipline as the Kelsenian-formalist disease the 1922 essay diagnosed,
as Part~\ref{sec:pattern} conceded.
The Article does not have a knockdown response to that reader.
Its position is that the formal wrapper is a constraint on the exercise of authority distinct from a
constraint on the decision to exercise it,
and that constraints on the exercise are valuable to the readership for whom such constraints are
in-principle available,
even if the strong Schmittian holds that constraints on the decision are unavailable on conceptual grounds.

The Article therefore concedes that its argument is addressed to a readership that does not share the strong
Schmittian view.
The readership it does address is the constitutional-law academy in its mainstream post-Weimar formation,
for which the legitimacy of formal constraint on emergency authority is a contested but live question
rather than a category mistake.
A reader who finds the Article's structural claim unconvincing because the reader holds the strong
Schmittian view is, on the Article's framing, holding a view that the Article does not purport to refute
and against which the Article's argument is not addressed.
The concession is offered not as evasion of the objection but as candor about the scope within which the
structural claim is meant to do work.

\subsection{The legal-realist objection}

A related but distinct objection comes from the legal-realist tradition.
On this view, the formal text of an authority is a poor predictor of how the authority will be exercised in practice.
The actors who exercise the authority will, predictably, develop interpretive practices that bend the text to
their operational needs.
A typed-rollback discipline imposed by statute would be, on the legal-realist view, subject to the same bending.
The witness $w$ would, in practice, become a pro forma artifact that satisfied the formal requirement without
imposing any meaningful constraint on the action it accompanies.

The legal-realist objection is substantial.
The empirical record on procedural requirements in administrative law and constitutional law is mixed.
Some procedural requirements have been highly effective in constraining the actions of officials
(notice and comment in administrative rulemaking),
and others have been substantially eroded through interpretive practice
(the warrant requirement in Fourth Amendment doctrine, where the catalog of exceptions has grown over time).
The grammar's typed-rollback witness might end up in either category.

The Article's response is that the grammar's discipline differs in kind from procedural requirements in
familiar legal-doctrinal frameworks.
A typed-rollback witness is not a procedure;
it is a construction that either type-checks or does not.
The witness cannot be satisfied by going through the motions of construction without producing a valid object.
The substrate's runtime kernel evaluates the witness mechanically and rejects ill-typed inputs.
A pro forma witness that did not type-check would fail admission at the construction step,
not at a later audit step.
The discipline is mechanical in a way that traditional procedural requirements are not.

This response has limits.
The discipline is mechanical at the level of the substrate.
It is not mechanical at the level of the substantive content of the witness.
A witness for a file quarantine that specifies a constructive path back to the file's original state is type-checked
at the substrate level;
whether the constructive path is in fact a meaningful restoration of the prior state is a substantive question
that the type system does not resolve.
A pro forma witness for a more abstract action -- restoring the geopolitical state of a country to its
pre-occupation condition, for example -- could in principle type-check while failing to provide any meaningful
reversibility.
The grammar provides a structural discipline at the wrapper level;
it does not, by itself, prevent substantive mischief at the content level.

\subsection{The doctrinal-fit objection}

A third objection is that the grammar's vocabulary -- typed witnesses, type-checking, refinement --
is alien to the legal-doctrinal vocabulary in which delegated emergency authority is debated.
A legal-academy reader may legitimately wonder whether the structural claim is a genuine addition to the doctrinal
discussion or a semantic relabeling of points that the existing doctrinal vocabulary already captures.
Sunset clauses, periodic reauthorization requirements, supermajority defaults for continuation,
and post-hoc review by independent bodies are all features of the existing doctrinal toolkit.
Does the typed-rollback discipline add anything that these existing features do not?

The Article's response is that the typed-rollback discipline adds something specific that the existing toolkit
does not directly provide:
the requirement is at the admission step rather than at the lapse step or the review step.
A sunset clause operates at the lapse step:
the authority lapses at a specified time unless extended.
The actions taken under the authority are not, by virtue of the sunset, reversed when the authority lapses.
A periodic reauthorization requirement operates at a review step:
the authority is reviewed periodically and may be continued or terminated.
The actions taken between reviews are not, by virtue of the review, reversed.
A supermajority default for continuation operates at the extension step:
the default is termination unless extended by supermajority.
Again, the actions taken under the authority are not reversed by the operation of the default.

The typed-rollback discipline differs from all three.
It operates at the admission step.
An action that cannot exhibit a constructive path back to the prior state cannot be admitted.
The discipline addresses the irreversibility of the action itself, not the lapse, the review, or the continuation of the authority.
This is, the Article argues, a genuine addition to the doctrinal toolkit rather than a relabeling of an existing feature.

The argument is not knockdown.
A legal-academy reader may reasonably observe that the doctrinal toolkit's mix of sunsets, reauthorizations,
and supermajority defaults addresses the irreversibility problem indirectly,
by making the authority itself less likely to persist long enough for irreversibility to matter.
The typed-rollback discipline addresses the problem directly.
Whether the direct address is preferable to the indirect address is a substantive question on which the Article
does not need to take a position;
the Article needs only the narrower claim that the direct address is structurally distinct from the indirect ones
and is not a relabeling.

\subsection{The cross-disciplinary asymmetry}

A final limit is the cross-disciplinary asymmetry of expertise between the Article's home discipline (formal methods,
applied to programmable governance) and the Article's target audience (legal academy, particularly constitutional law).
The author is not a lawyer.
The Article relies on legal historiography, doctrinal accounts, and case-study literatures that a constitutional-law
scholar would treat with greater facility than a non-lawyer can.
The Article has tried, in each case study and in the doctrinal framing, to defer to the secondary literature where
the Article's claim depends on a doctrinal characterization,
and to flag uncertainty where the doctrinal record is contested.

The cross-disciplinary asymmetry is, on the author's view, a reason for the Article to be co-authored or extensively
reviewed by a constitutional-law scholar.
The structural claim does not depend on getting any particular doctrinal claim exactly right;
it depends on getting the structural pattern across the five case studies correct,
which is a claim that survives reasonable variation in the doctrinal characterizations of any individual case.
A reader who finds particular doctrinal claims in Part~\ref{sec:cases} unconvincing is invited to test the structural
claim against the case studies that the reader finds most defensible.
The grammar's claim is robust to the failure of any one or two case studies;
the claim depends on the pattern across the five.
